%====================================
% KAPITEL 2: THEORETISCHE GRUNDLAGEN
%====================================
\section{Theoretische Grundlagen}

Das folgende Kapitel legt das begriffliche und theoretische Fundament der Arbeit. Zun\"achst werden ERP-Systeme definiert, in ihrer historischen Entwicklung eingeordnet und die Besonderheiten ihres Einsatzes in mittelst\"andischen Unternehmen herausgearbeitet (Kapitel~2.1). Anschlie\ss{}end wird mit dem IS-Success-Modell nach DeLone und McLean die Erfolgsperspektive der Untersuchung eingef\"uhrt (Kapitel~2.2), bevor das ADKAR-Modell die Change-Management-Perspektive erg\"anzt (Kapitel~2.3). Kapitel~2.4 verdichtet beide Perspektiven zu einem vorl\"aufigen Bezugsrahmen, der die systematische Literaturanalyse in den Kapiteln~3 und~4 strukturiert.

\subsection{Enterprise-Resource-Planning-Systeme und ihre Besonderheiten im Mittelstand}

\textbf{Begriff und Merkmale von ERP-Systemen}

Enterprise-Resource-Planning-Systeme (ERP-Systeme) sind integrierte, modular aufgebaute Standardsoftwaresysteme, die die zentralen Gesch\"aftsprozesse eines Unternehmens -- etwa Finanz- und Rechnungswesen, Materialwirtschaft, Produktion, Vertrieb und Personal -- funktions\"ubergreifend auf Basis einer gemeinsamen Datenbank abbilden und steuern.\footcite[Vgl.][S.~30--36]{holland1999} Das konstituierende Merkmal ist dabei die Integration: Anstelle isolierter Insell\"osungen f\"ur einzelne Funktionsbereiche werden s\"amtliche betriebswirtschaftlich relevanten Daten in einem konsistenten Datenmodell zusammengef\"uhrt, sodass Informationen bereichs\"ubergreifend in Echtzeit verf\"ugbar sind.\footcite[Vgl.][S.~30--36]{holland1999} Davenport charakterisierte ERP-Systeme bereits fr\"uh als Enterprise Systems, deren Einf\"uhrung weit \"uber ein reines IT-Projekt hinausgeht, da die Software den Gesch\"aftsprozessen eine eigene Logik aufpr\"agt und Unternehmen ihre Abl\"aufe h\"aufig an die Referenzprozesse der Software anpassen m\"ussen -- und nicht umgekehrt.\footcite[Vgl.][S.~121--131]{davenport1998}

Historisch haben sich ERP-Systeme aus den Material-Requirements-Planning-Systemen (MRP) der 1970er-Jahre und deren Erweiterung zum Manufacturing Resource Planning (MRP~II) in den 1980er-Jahren entwickelt. Mit der Ausdehnung des Funktionsumfangs auf s\"amtliche Unternehmensressourcen etablierte sich Anfang der 1990er-Jahre der Begriff des Enterprise Resource Planning.\footcite[Vgl.][S.~258--259]{somers2004} J\"ungere Entwicklungen erweitern das klassische ERP-Verst\"andnis: Cloud-basierte Bereitstellungsmodelle (Software as a Service) senken die anf\"anglichen Investitionskosten und den internen Betriebsaufwand, verlagern jedoch Anpassungs- und Integrationsfragen in die Konfiguration und erh\"ohen die Abh\"angigkeit vom Anbieter.\footcite[Vgl.][S.~3--4]{butarbutar2023} Gerade f\"ur ressourcenschw\"achere Unternehmen gewinnen solche Modelle an Bedeutung, was den ab 2015 gew\"ahlten Betrachtungszeitraum dieser Arbeit zus\"atzlich begr\"undet.

Aus prozessualer Sicht l\"asst sich die ERP-Einf\"uhrung als mehrphasiger Lebenszyklus beschreiben, der von der Auswahl- und Planungsphase \"uber die eigentliche Implementierung bis zur Stabilisierungs- und Weiterentwicklungsphase reicht.\footcite[Vgl.][S.~258--259]{somers2004} F\"ur die vorliegende Arbeit ist diese Perspektive relevant, weil Change-Management-Ma\ss{}nahmen in den einzelnen Phasen unterschiedlich wirken -- etwa Kommunikation in der Ank\"undigungsphase, Schulung unmittelbar vor und nach dem Go-Live.

\textbf{Abgrenzung des Mittelstandsbegriffs}

Mittelst\"andische Unternehmen unterscheiden sich in mehreren Dimensionen grundlegend von Gro\ss{}unternehmen. Sie zeichnen sich typischerweise durch flachere Hierarchien, k\"urzere Entscheidungswege, eine h\"ohere Bedeutung von Schl\"usselpersonen sowie begrenzte finanzielle und personelle Ressourcen aus.\footcite[Vgl.][S.~3--4]{butarbutar2023} Diese Spezifika haben unmittelbare Auswirkungen auf ERP-Implementierungsprojekte: W\"ahrend Gro\ss{}konzerne auf umfangreiche interne IT-Abteilungen und spezialisierte Projektteams zur\"uckgreifen k\"onnen, sind mittelst\"andische Unternehmen h\"aufig auf externe Berater angewiesen und verf\"ugen "uber weniger personelle Kapazit\"aten f\"ur das begleitende Change Management. Gleichzeitig erm\"oglichen die k\"urzeren Entscheidungswege eine h\"ohere Umsetzungsgeschwindigkeit, sofern die Gesch\"aftsf\"uhrung das Projekt aktiv vorantreibt.

Die Komplexit\"at von ERP-Einf\"uhrungen wird durch die Notwendigkeit verst\"arkt, bestehende Gesch\"aftsprozesse an die Anforderungen der Software anzupassen (Standardisierung) oder die Software an die unternehmensspezifischen Prozesse anzupassen (Individualisierung).\footcite[Vgl.][S.~258--259]{somers2004} Insbesondere der Mittelstand steht dabei vor der Herausforderung, zwischen branchen\"ublichen Best Practices und eigenen, oft historisch gewachsenen Prozessstrukturen abzuw\"agen.
